\section{Ablations and Analysis}
\label{sec:ablation}

% Rows and cells marked "--" are planned experiments without results yet.
% Replace them, and rewrite the matching "Planned" paragraphs, once runs finish.

\subsection{The cost unit decides the outcome}
\label{sec:unitresult}

To test the choice of cost unit (Section~\ref{sec:problem}), we train the parity rule with
cost measured per byte instead of per parallel sentence (Table~\ref{tab:byte}). Per byte,
English becomes the worst-off language, and the parity rule gives it 99\% of the merges.
Measured by what a user pays, the result is the opposite of the intention: the Indian
languages need about twice as many tokens per sentence as under classical BPE (63.6 vs.\
35.0 for Hindi), while English needs fewer (26.6 vs.\ 29.0). Thus, a per-byte cost inverts
the parity rule across scripts, and the cost must be measured on shared content.

\begin{table}[t]
\centering
\small
\resizebox{\columnwidth}{!}{%
\begin{tabular}{lcccc}
\toprule
Tokenizer & Hindi & Marathi & Gujarati & English\\
\midrule
Classical BPE & 35.0 & 32.5 & 34.4 & 29.0\\
Parity-aware, cost per byte & 63.6 & 70.6 & 66.5 & 26.6\\
Parity-aware, cost per sentence & 32.1 & 32.5 & 32.3 & 32.5\\
\bottomrule
\end{tabular}}
\caption{Tokens per sentence on devtest when the training cost is measured per byte
(single run) or per parallel sentence (five-sample mean, Table~\ref{tab:perlang}).}
\label{tab:byte}
\end{table}

\subsection{Data balance and vocabulary size}
\label{sec:sens}

To measure how much the disparity depends on the setting, we repeat the experiment with
other amounts of English and with larger vocabularies (Table~\ref{tab:sens}). The direction
is the same in every setting: the Gini coefficient and the compression fall at every step
in $p$. However, the size of the problem changes considerably with data balance. With equal
amounts of text, classical BPE is already close to parity (Gini 0.019), and full parity costs
1.1\% of compression. With eight times as much English, the Gini coefficient is 0.051, a
Hindi sentence costs 30\% more tokens than its English translation, and full parity costs
7.0\%. Moreover, with equal data the remaining disparity is not between English and the
Indian languages, since English (31.7 tokens) lies between Marathi (30.4) and Hindi (33.3).

Doubling the vocabulary to 64k leaves the Gini coefficient of both classical BPE (0.038) and
Parity-aware BPE (0.003) unchanged. Thus, between 32k and 64k we do not observe the loss of
equity that \citet{lee2026} report for Parity-aware BPE at larger vocabularies.

\noindent\textbf{Planned: 128k vocabulary.}
Since Parity-aware BPE bounds only the worst-off language, \citet{lee2026} find that the
dominant language keeps acquiring merges as the vocabulary grows. For finite $p$, every
language keeps a positive weight in Eq.~\eqref{eq:rule}, so we expect intermediate $p$ to
be less affected. At 128k, the Gini coefficient is -- for classical BPE, -- for $p=8$ and --
for Parity-aware BPE.

\begin{table*}[t]
\centering
\small
\begin{tabular}{llcccccc}
\toprule
English & Vocab. & $p=1$ & $p=2$ & $p=4$ & $p=8$ & $p=\infty$ & Cost of parity\\
\midrule
$1\times$ & 32k & 0.0186 & 0.0172 & 0.0148 & 0.0117 & 0.0034 & 1.1\%\\
$4\times$ $\dagger$ & 32k & 0.0380 & 0.0347 & 0.0298 & 0.0229 & 0.0028 & 4.0\%\\
$8\times$ & 32k & 0.0505 & 0.0461 & 0.0383 & 0.0298 & 0.0033 & 7.0\%\\
$4\times$ & 64k & 0.0378 & 0.0353 & 0.0313 & 0.0255 & 0.0032 & 5.1\%\\
$4\times$ & 128k & -- & -- & -- & -- & -- & --\\
\bottomrule
\end{tabular}
\caption{Gini coefficient of per-language cost in other settings, and the compression lost
by going from classical to Parity-aware BPE. The row marked $\dagger$ is the five-sample mean
of Table~\ref{tab:front}; the others are single runs.}
\label{tab:sens}
\end{table*}

\subsection{Comparison with Worst-$N$ parity (planned)}
\label{sec:worstn}

The Worst-$N$ rule \citep{kanjirangat2026} offers a second way to interpolate, by
optimising the $N$ worst-off languages jointly. To compare it with $p$-fair BPE, we train
$N\in\{1,2,3\}$ in the same setting (Table~\ref{tab:worstn}); $N=1$ coincides with the
parity rule. At matched compression, the Gini difference between Worst-$N$ and $p$-fair BPE
is --.

\begin{table}[t]
\centering
\small
\begin{tabular}{lcc}
\toprule
Tokenizer & Gini $\downarrow$ & Compression $\uparrow$\\
\midrule
Worst-$N$, $N=1$ ($=p{=}\infty$) & 0.0028 & 6.634\\
Worst-$N$, $N=2$ & -- & --\\
Worst-$N$, $N=3$ & -- & --\\
$p$-fair, $p=8$ & 0.0229 & 6.868\\
$p$-fair, $p=32$ & 0.0084 & 6.760\\
\bottomrule
\end{tabular}
\caption{Worst-$N$ parity against $p$-fair BPE at 32k with English at $4\times$.
Compression is bytes per token on the training text.}
\label{tab:worstn}
\end{table}

\subsection{Morphology penalty (planned)}
\label{sec:morph}

To add morphological alignment as a third axis, we subtract a penalty for merges that join
symbols from different morphemes:
\begin{equation}
  (a,b)^\star = \argmax_{(a,b)}\; s_p(a,b) - \lambda\,\beta(a,b),
  \label{eq:morph}
\end{equation}
where $\beta(a,b)$ counts the occurrences of the pair that cross a morpheme boundary in the
SuTRA segmentations \citep{rathore2026}, and $\lambda\ge0$ sets the strength. Since both the
count and the penalty of a pair fall over time, the score is no longer guaranteed to be
non-increasing, and Proposition~\ref{prop:lazy} requires a modified bound for $\lambda>0$.
Table~\ref{tab:morph} reports boundary F1 against the SuTRA gold segmentations. Increasing
$\lambda$ changes boundary F1 by --, at a cost of -- in Gini and -- in compression.

\begin{table}[t]
\centering
\small
\resizebox{\columnwidth}{!}{%
\begin{tabular}{llccccc}
\toprule
 & & & & \multicolumn{3}{c}{Boundary F1 $\uparrow$}\\
\cmidrule(lr){5-7}
$p$ & $\lambda$ & Gini $\downarrow$ & Compr. $\uparrow$ & hi & mr & gu\\
\midrule
1 & 0 & 0.0380 & 6.910 & -- & -- & --\\
1 & -- & -- & -- & -- & -- & --\\
8 & 0 & 0.0229 & 6.868 & -- & -- & --\\
8 & -- & -- & -- & -- & -- & --\\
$\infty$ & 0 & 0.0028 & 6.634 & -- & -- & --\\
$\infty$ & -- & -- & -- & -- & -- & --\\
\bottomrule
\end{tabular}}
\caption{Effect of the morphology penalty. Compression is bytes per token on the training
text; boundary F1 is measured against the SuTRA gold segmentations.}
\label{tab:morph}
\end{table}

\subsection{Vocabulary allocation across scripts (planned)}
\label{sec:alloc}

A complementary question is how a fixed vocabulary of size $K$ should be divided between
scripts (Devanagari, Gujarati and Latin). We fit a cost curve $t_c(k)=a_c k^{-b_c}+d_c$ per
script $c$ and solve $\min_k \sum_c \pi_c\, t_c(k_c)^p$ subject to $\sum_c k_c=K$, which is
convex for $p\ge1$ and is solved by water-filling (Appendix~\ref{app:alloc}). We compare it
with a uniform split, the union-proportion rule of \citet{brahma2026} and VoCap
\citep{zheng2021vocap} (Table~\ref{tab:alloc}).

\begin{table}[t]
\centering
\small
\begin{tabular}{lcc}
\toprule
Allocation & Gini $\downarrow$ & Compression $\uparrow$\\
\midrule
Uniform split & -- & --\\
Union proportion \citep{brahma2026} & -- & --\\
VoCap \citep{zheng2021vocap} & -- & --\\
Convex, $p=1$ & -- & --\\
Convex, $p=8$ & -- & --\\
\bottomrule
\end{tabular}
\caption{Vocabulary allocation across scripts at a total budget of 32k.}
\label{tab:alloc}
\end{table}

\subsection{Downstream evaluation (planned)}
\label{sec:downstream}

To test whether intrinsic gains matter in use, we evaluate Hindi--Marathi translation and
the number of tokens per technical term (Table~\ref{tab:downstream}).

\begin{table}[t]
\centering
\small
\resizebox{\columnwidth}{!}{%
\begin{tabular}{lccc}
\toprule
Tokenizer & chrF++ hi$\to$mr $\uparrow$ & chrF++ mr$\to$hi $\uparrow$ & Tok./term $\downarrow$\\
\midrule
Classical ($p=1$) & -- & -- & --\\
$p$-fair, $p=8$ & -- & -- & --\\
Hybrid, 75\% classical & -- & -- & --\\
Parity-aware ($p=\infty$) & -- & -- & --\\
\bottomrule
\end{tabular}}
\caption{Downstream evaluation.}
\label{tab:downstream}
\end{table}